\section{Limitações}\label{sec:limitacoes}

A curadoria removeu as \svAusentes{} linhas com valor ausente. Elas ganham mais de 50 mil com menos frequência (\svPositivosAusentes\% contra \svPositivos\%), de modo que a base ficou um pouco mais ``rica'' e os modelos não aprenderam o perfil de quem está fora do mercado de trabalho ou tem vínculo informal. Tratar o ausente como uma categoria própria manteria essas pessoas; não medimos quanto isso mudaria os resultados.

As grades da SVM foram escolhidas a partir de um estudo numa única divisão, a seed 0, que também faz parte das 10 seeds da comparação; o mapa de $C$ e $\gamma$ é largo e plano o bastante para que isso pese pouco, mas a escolha não foi independente da avaliação. A busca usou a acurácia como critério em todos os modelos, e a Seção~\ref{sec:sv_seeds} mostra que esse critério piora a AUC da SVM RBF; com a AUC como critério, a comparação poderia mudar. A Random Forest foi usada na configuração padrão, sem ajuste. As 10 seeds variam o sorteio da divisão, mas os dados são os mesmos, de modo que o desvio padrão não mede a incerteza para dados novos. Os tempos foram medidos numa única máquina, num núcleo só; os valores absolutos mudam com o computador, e as razões entre modelos são mais estáveis que os tempos em si.

Deixar raça e sexo fora dos atributos não impede que o modelo aprenda padrões ligados a eles: outros atributos carregam a mesma informação, como o relacionamento, em que ``marido'' e ``esposa'' indicam o sexo quase sem exceção. Não avaliamos aqui os efeitos disso. A base é de 1994, e o limite de 50 mil dólares, de outro tempo; as conclusões valem para esta base.
